% GENERATED FILE. Edit canonical lessons, then run npm run generate:books.
% canonical-source-hash: fnv1a64:05477d145c95356d
% canonical-lessons: AR-C35-jar, AR-C35-ibn-bint, AR-C35-sadiq

\chapter{The People You Introduce}
\label{ch:ar-people-you-introduce}
\hlchaptermodality{35}

\begin{chapteropening}
\textbf{By the end of this chapter:} \emph{I can name and introduce a neighbour, a son, a daughter and a friend in Arabic, and say what single letter genders a pair of them.}

Take ṣadīq back to \ar{ص}-\ar{د}-\ar{ق}, 'to be truthful', and run the family ṣidq, ṣādiq, ṣadaqa and aṣdiqāʾ, naming the faʿīl pattern it shares with ḥalīb and malīḥ and the doer shape it shares with qāʾil and kātib, then line up the people already owned — ab, umm, akh, ukht — against jārī, ibn, bint and ṣadīqī.
\end{chapteropening}

\section[jār]{\ar{جار} (jār) — \textquotedblleft{}neighbour,\textquotedblright{} and a root with a letter in hiding}

\label{lesson:AR-C35-jar}



\begin{quote}
After a chapter of things you order, a chapter of people you introduce — starting with the one who lives next door.
\end{quote}

\subsection*{The letters in this word}
Nothing new, and it is short: \textbf{\ar{ج}} (\emph{jīm}, the hook-and-tail body with one dot below), \textbf{\ar{ا}} (\emph{alif}) and \textbf{\ar{ر}} (\emph{rā}). Both \textbf{\ar{ا}} and \textbf{\ar{ر}} are non-joiners, so \textbf{\ar{جار}} barely joins at all — compare \textbf{\ar{طعام}}, which broke once before its \emph{mīm}.

\begin{cousinweb}
\textbf{\ar{جار}} (\emph{jār}) = \textquotedblleft{}\textbf{a neighbour}.\textquotedblright{} Add the \textquotedblleft{}my\textquotedblright{} ending of \textbf{\ar{اسمي}} (\emph{ismī}) and you have \textbf{\ar{جاري}} (\emph{jārī}), \textquotedblleft{}\textbf{my neighbour}\textquotedblright{} — which is how the word usually arrives in a conversation.

Now listen to the root. You can hear only \emph{j} and \emph{r}. But an Arabic root wants three consonants, and the missing one is in hiding: the root is \textbf{\ar{ج}-\ar{و}-\ar{ر}}, and its \textbf{\ar{و}} has collapsed into that long \textbf{\ar{ا}}.

You have watched this before. \textbf{\ar{قال}} (\emph{qāla}), \textquotedblleft{}he said,\textquotedblright{} is \textbf{\ar{ق}-\ar{و}-\ar{ل}} with the same disappearing \textbf{\ar{و}}. Arabic calls these \textbf{hollow} roots, and the hidden letter comes back when the word changes shape:

\begin{itemize}
\raggedright
  \item \textbf{\ar{جاوَرَ}} (\emph{jāwara}) — \textquotedblleft{}he lived alongside.\textquotedblright{} There is the \textbf{\ar{و}}, in the open.
  \item \textbf{\ar{جيران}} (\emph{jīrān}) — \textquotedblleft{}\textbf{neighbours},\textquotedblright{} where it returns as a long \textbf{ī}.
\end{itemize}

Hebrew's \textbf{\he{גֵּר}} (\emph{ger}), a resident outsider living among a people, comes from the same Semitic root for dwelling alongside.

English \textbf{neighbour} is no relative — it is Germanic, the \emph{nigh-dweller}. Two unrelated languages built one word out of one idea, worth noticing \textbf{because} it is not descent.
\end{cousinweb}

\subsection*{Your turn}
Say these aloud:

\begin{itemize}
\raggedright
  \item jīm, alif, rā — \textquotedblleft{}jār,\textquotedblright{} a neighbour
  \item with the my ending — \textquotedblleft{}jārī,\textquotedblright{} my neighbour, as ismī is my name
  \item the hidden letter — the root is j-w-r, and the wāw is inside the alif
  \item the same hiding place — qāla, he said, from q-w-l
  \item where it comes back — jāwara; jīrān, neighbours
  \item the hook family by dots — ḥāʾ, khāʾ, jīm
  \item the restaurant you built — maṭʿam, the place of eating
  \item the honest non-cousin — English neighbour, the nigh-dweller
\end{itemize}

\begin{tcolorbox}[breakable,colback=teal!4,colframe=teal!35!black,title={Before you move on}]
What is \textbf{\ar{جار}}'s third root letter, and where has it gone? (\textbf{\ar{و}} — collapsed into the long \emph{alif}, as in \emph{\textbf{qāla}} from \textbf{\ar{ق}-\ar{و}-\ar{ل}}.) Where does it resurface? (\textbf{In \emph{jāwara} and \emph{jīrān}.}) How do you say \textquotedblleft{}my neighbour\textquotedblright{}? (\emph{\textbf{Jārī}} — the ending of \emph{ismī}.) Is English \emph{neighbour} related? (\textbf{No} — Germanic, the \emph{nigh-dweller}.)
\end{tcolorbox}
\section[ibn, bint]{\ar{ابن}, \ar{بنت} (ibn, bint) — son and daughter, one letter apart}

\label{lesson:AR-C35-ibn-bint}



\begin{quote}
A pair you can learn in one move, because Arabic has already shown you the move it uses.
\end{quote}

\subsection*{The letters in this word}
Nothing new in either. \textbf{\ar{ابن}} is \textbf{\ar{ا}} (\emph{alif}), \textbf{\ar{ب}} (\emph{bāʾ}, one dot below) and \textbf{\ar{ن}} (\emph{nūn}, one dot above) — two letters that share a bowl and differ only by where the dot sits.

\textbf{\ar{بنت}} is the same \textbf{\ar{ب}} and \textbf{\ar{ن}}, then \textbf{\ar{ت}} (\emph{tāʾ}, two dots above) on the end. Both words join throughout, unlike \textbf{\ar{جار}}, which stopped twice.

\begin{cousinweb}
\textbf{\ar{ابن}} (\emph{ibn}) = \textquotedblleft{}\textbf{a son}.\textquotedblright{} \textbf{\ar{بنت}} (\emph{bint}) = \textquotedblleft{}\textbf{a daughter}.\textquotedblright{}

What separates them is \textbf{one \ar{ت}}, and nothing else.

You have met that letter doing that job. \textbf{\ar{أخ}} (\emph{akh}) is a brother; \textbf{\ar{أخت}} (\emph{ukht}) is a sister. Same word, one \textbf{\ar{ت}} added for the feminine. Arabic does not build a separate word for the woman; it marks the word it has. (This is the bare cousin of the \textbf{\ar{ة}} that marks \emph{qahwa} feminine — same letter, tied instead of open.)

On the root, tradition and comparison disagree. Arab grammarians file \textbf{\ar{ابن}} under \textbf{\ar{ب}-\ar{ن}-\ar{ي}}, \textquotedblleft{}\textbf{to build}\textquotedblright{} — a son as what a house is built of, which is a lovely idea. Comparative scholars are less sure: Hebrew \textbf{\he{בֵּן}} (\emph{ben}) and \textbf{\he{בַּת}} (\emph{bat}) show the same bare shape, more like a very old primitive word than something derived from a verb. Hold it loosely.

Hebrew \emph{bat} is \emph{bint} with the \textbf{n} swallowed into the \textbf{t}. And \emph{ibn} stands in plain sight inside names English carries: \textbf{Ibn Sīnā}, the physician the Latin west called Avicenna. No English \textbf{word} descends from it.
\end{cousinweb}

\subsection*{Your turn}
Say these aloud:

\begin{itemize}
\raggedright
  \item alif, bāʾ, nūn — \textquotedblleft{}ibn,\textquotedblright{} a son
  \item bāʾ, nūn, tāʾ — \textquotedblleft{}bint,\textquotedblright{} a daughter
  \item the dots that separate bāʾ and nūn — one below, one above
  \item the move you already knew — akh, then ukht
  \item and the parents beside them — ab, a father; umm, a mother
  \item the loosely held root — b-n-y, to build, as the grammarians have it
  \item the Hebrew pair — ben, bat
  \item the neighbour, and the emphatic ṭāʾ of ṭaʿām — jārī; ṭaʿām, maṭʿam
\end{itemize}

\begin{tcolorbox}[breakable,colback=teal!4,colframe=teal!35!black,title={Before you move on}]
What is the only difference between \textbf{\ar{ابن}} and \textbf{\ar{بنت}}? (\textbf{An added \ar{ت}} — as with \emph{\textbf{akh}} and \emph{\textbf{ukht}}.) And the parents? (\emph{\textbf{Ab}} and \emph{\textbf{umm}}.) What root do grammarians give \emph{ibn}, and why hold it loosely? (\textbf{\ar{ب}-\ar{ن}-\ar{ي}, \textquotedblleft{}to build\textquotedblright{}} — but Hebrew \emph{ben} and \emph{bat} suggest an older bare word.) Where does \emph{ibn} appear in a name English knows? (\emph{\textbf{Ibn Sīnā}}, Avicenna.)
\end{tcolorbox}
\section[ṣadīq]{\ar{صديق} (ṣadīq) — \textquotedblleft{}friend,\textquotedblright{} the one who is true to you}

\label{lesson:AR-C35-sadiq}



\begin{quote}
The chapter closes on a word whose root is not about company at all. It is about telling the truth.
\end{quote}

\subsection*{The letters in this word}
Nothing new. \textbf{\ar{ص}} (\emph{ṣād}) is the emphatic \textbf{\ar{س}}, first met in \textbf{\ar{صباح}} (\emph{ṣabāḥ}) — heavy and dark where \emph{sīn} is light. Then \textbf{\ar{د}} (\emph{dāl}), \textbf{\ar{ي}} (\emph{yāʾ}) carrying the long \textbf{ī}, and \textbf{\ar{ق}} (\emph{qāf}). \textbf{\ar{د}} is a non-joiner, so \textbf{\ar{صديق}} breaks after it: \textbf{\ar{صد}}, then \textbf{\ar{يق}}.

The feminine is \textbf{\ar{صَديقة}} (\emph{ṣadīqa}) — the \textbf{\ar{ة}}, tied on the end, doing what it did to \emph{qahwa}: the open \textbf{\ar{ت}} of \textbf{\ar{بنت}} and \textbf{\ar{أخت}} in its tied form.

\begin{cousinweb}
\textbf{\ar{صَديق}} (\emph{ṣadīq}) = \textquotedblleft{}\textbf{a friend}.\textquotedblright{} Root \textbf{\ar{ص}-\ar{د}-\ar{ق}}, and the root means \textquotedblleft{}\textbf{to be truthful, to be sincere}.\textquotedblright{}

A friend, in Arabic, is defined by honesty rather than proximity. The family agrees:

\begin{itemize}
\raggedright
  \item \textbf{\ar{صِدْق}} (\emph{ṣidq}) — \textquotedblleft{}\textbf{truth}, sincerity.\textquotedblright{}
  \item \textbf{\ar{صادِق}} (\emph{ṣādiq}) — \textquotedblleft{}\textbf{truthful},\textquotedblright{} the \textbf{doer} shape worn by \textbf{\ar{قائِل}} (\emph{qāʾil}, a sayer) and \textbf{\ar{كاتِب}} (\emph{kātib}, a writer).
  \item \textbf{\ar{صَدَقة}} (\emph{ṣadaqa}) — \textquotedblleft{}\textbf{voluntary alms}.\textquotedblright{} A gift given for its own sake: a \textbf{true} gift, as against one owed.
  \item \textbf{\ar{أَصْدِقاء}} (\emph{aṣdiqāʾ}) — \textquotedblleft{}\textbf{friends}.\textquotedblright{}
\end{itemize}

And now the pattern. \textbf{\ar{صَديق}} is \textbf{\ar{فَعيل}} (\emph{faʿīl}) — the shape of \textbf{\ar{حَليب}} (\emph{ḥalīb}), the milked thing, and \textbf{\ar{مَليح}} (\emph{malīḥ}), the salty one. Root plus \emph{faʿīl} hands back a quality made into a person or a thing.

This is the machinery that turned \textbf{\ar{سَلام}} (\emph{salām}) into \textbf{\ar{السَّلامة}} (\emph{as-salāma}): one root, several patterns, several words.

English took nothing from \textbf{\ar{ص}-\ar{د}-\ar{ق}}. \emph{Friend} is Germanic, from an old word for loving.
\end{cousinweb}

\subsection*{Your turn}
Say these aloud:

\begin{itemize}
\raggedright
  \item ṣād, dāl, yāʾ, qāf — \textquotedblleft{}ṣadīq,\textquotedblright{} a friend
  \item the emphatic pair — sīn light, ṣād heavy, as in ṣabāḥ
  \item what the root means — to be truthful; a friend is one who is true
  \item the family — ṣidq, ṣādiq, ṣadaqa, aṣdiqāʾ
  \item the doer shape — qāʾil, kātib, ṣādiq
  \item the faʿīl shape — ḥalīb, malīḥ, ṣadīq
  \item the feminine — ṣadīqa, tied tāʾ, as as-salāma has it
  \item this chapter's people — jārī with its two non-joiners, ibn, bint, ṣadīqī
  \item and the family you had — ab, umm, akh, ukht
\end{itemize}

\begin{tcolorbox}[breakable,colback=teal!4,colframe=teal!35!black,title={Before you move on}]
What does \textbf{\ar{ص}-\ar{د}-\ar{ق}} mean, and what does that make a friend? (\textbf{To be truthful} — the one who is true to you.) Name the family. (\emph{\textbf{Ṣidq}}, \emph{\textbf{ṣādiq}}, \emph{\textbf{ṣadaqa}}, \emph{\textbf{aṣdiqāʾ}}.) Which earlier words wear \emph{ṣadīq}'s \emph{faʿīl} pattern, and which wear \emph{ṣādiq}'s doer shape? (\emph{\textbf{Ḥalīb}} and \emph{\textbf{malīḥ}}; \emph{\textbf{qāʾil}} and \emph{\textbf{kātib}}.) Introduce the people you can now name. (\emph{\textbf{Jārī}}, \emph{\textbf{ibn}}, \emph{\textbf{bint}}, \emph{\textbf{ṣadīqī}} — beside \emph{\textbf{ab}}, \emph{\textbf{umm}}, \emph{\textbf{akh}}, \emph{\textbf{ukht}}.) What does the \textbf{\ar{ة}} of \emph{ṣadīqa} do? (\textbf{Marks it feminine} — as on \emph{qahwa} and \emph{as-salāma}.)
\end{tcolorbox}
